
\documentclass{article}
\usepackage{colortbl}
\usepackage{makecell}
\usepackage{multirow}
\usepackage{supertabular}

\begin{document}

\newcounter{utterance}

\centering \large Interaction Transcript for game `dond', experiment `semi\_it', episode 16 with gpt{-}5{-}2025{-}08{-}07{-}no{-}reasoning{-}t1.0.
\vspace{24pt}

{ \footnotesize  \setcounter{utterance}{1}
\setlength{\tabcolsep}{0pt}
\begin{supertabular}{c@{$\;$}|p{.15\linewidth}@{}p{.15\linewidth}p{.15\linewidth}p{.15\linewidth}p{.15\linewidth}p{.15\linewidth}}
    \# & $\;$A & \multicolumn{4}{c}{Game Master} & $\;\:$B\\
    \hline

    \theutterance \stepcounter{utterance}  
    & & \multicolumn{4}{p{0.6\linewidth}}{
        \cellcolor[rgb]{0.9,0.9,0.9}{
            \makecell[{{p{\linewidth}}}]{
                \texttt{\tiny{[P1$\langle$GM]}}
                \texttt{State giocando a un gioco di negoziazione in cui dovete accordarvi su come dividere una serie di oggetti tra voi e un altro giocatore.} \\
\\ 
\texttt{Regole:} \\
\texttt{(a) A Lei e all'altro giocatore viene dato un insieme di oggetti. Ognuno di voi riceve anche una funzione di valore segreta, che rappresenta il valore di ciascun tipo di oggetto.} \\
\texttt{(b) Si scambiano messaggi con l'altro giocatore per concordare chi si aggiudica gli oggetti. Potete inviare un massimo di 5 messaggi ciascuno, oppure terminare in anticipo facendo una proposta segreta in qualsiasi momento.} \\
\texttt{(c) A ciascuno di voi viene chiesto di inviare una proposta segreta indicando gli oggetti che desiderate, formattata tra parentesi quadre come segue: "{[}Proposta: <numero> <nome oggetto>, <numero> <nome oggetto>, <...>{]}".} \\
\texttt{(d) Se le vostre proposte sono complementari, cioè ci sono abbastanza oggetti per soddisfare entrambe le proposte, a ciascun giocatore viene assegnato un punteggio basato sulla somma dei valori degli oggetti ricevuti. In caso contrario, entrambi ricevono zero punti.} \\
\texttt{(e) L'obiettivo è massimizzare il proprio punteggio.} \\
\\ 
\texttt{Cominciamo.} \\
\\ 
\texttt{L'insieme degli oggetti disponibili è:} \\
\\ 
\texttt{3 vasi, 1 sedia, 1 cavo} \\
\\ 
\texttt{La funzione valore segreta è:} \\
\\ 
\texttt{vaso: 2, sedia: 3, cavo: 1} \\
\\ 
\texttt{IMPORTANTE: i vostri messaggi, a meno che non si tratti di una proposta segreta, vengono trasmessi direttamente all'altro giocatore, quindi non includete alcuna risposta alle regole o testo di annuncio del vostro messaggio. Per fare una proposta segreta, utilizzate il formato indicato. Non utilizzare delle parentesi quadre quando si comunica all'altro giocatore, altrimenti verrà interpretata come una proposta segreta.} \\
            }
        }
    }
    & \\ \\

    \theutterance \stepcounter{utterance}  
    & \multicolumn{4}{p{0.6\linewidth}}{
        \cellcolor[rgb]{0.9,0.9,0.9}{
            \makecell[{{p{\linewidth}}}]{
                \texttt{\tiny{[P1$\rangle$GM]}}
                \texttt{Ciao! Per me i vasi valgono 2 ciascuno, la sedia 3 e il cavo 1. Ti propongo questo: io prendo la sedia e 1 vaso, tu prendi 2 vasi e il cavo. Ti va?} \\
            }
        }
    }
    & & \\ \\

    \theutterance \stepcounter{utterance}  
    & & & \multicolumn{2}{p{0.3\linewidth}}{
        \cellcolor[rgb]{0.9,0.9,0.9}{
            \makecell[{{p{\linewidth}}}]{
                \texttt{\tiny{[GM$|$GM]}}
                \texttt{continue} \\
            }
        }
    }
    & & \\ \\

    \theutterance \stepcounter{utterance}  
    & & & \multicolumn{2}{p{0.3\linewidth}}{
        \cellcolor[rgb]{0.9,0.9,0.9}{
            \makecell[{{p{\linewidth}}}]{
                \texttt{\tiny{[GM$|$GM]}}
                \texttt{Ciao! Per me i vasi valgono 2 ciascuno, la sedia 3 e il cavo 1. Ti propongo questo: io prendo la sedia e 1 vaso, tu prendi 2 vasi e il cavo. Ti va?} \\
            }
        }
    }
    & & \\ \\

    \theutterance \stepcounter{utterance}  
    & & \multicolumn{4}{p{0.6\linewidth}}{
        \cellcolor[rgb]{0.9,0.9,0.9}{
            \makecell[{{p{\linewidth}}}]{
                \texttt{\tiny{[GM$\rangle$P2]}}
                \texttt{State giocando a un gioco di negoziazione in cui dovete accordarvi su come dividere una serie di oggetti tra voi e un altro giocatore.} \\
\\ 
\texttt{Regole:} \\
\texttt{(a) A Lei e all'altro giocatore viene dato un insieme di oggetti. Ognuno di voi riceve anche una funzione di valore segreta, che rappresenta il valore di ciascun tipo di oggetto.} \\
\texttt{(b) Si scambiano messaggi con l'altro giocatore per concordare chi si aggiudica gli oggetti. Potete inviare un massimo di 5 messaggi ciascuno, oppure terminare in anticipo facendo una proposta segreta in qualsiasi momento.} \\
\texttt{(c) A ciascuno di voi viene chiesto di inviare una proposta segreta indicando gli oggetti che desiderate, formattata tra parentesi quadre come segue: "{[}Proposta: <numero> <nome oggetto>, <numero> <nome oggetto>, <...>{]}".} \\
\texttt{(d) Se le vostre proposte sono complementari, cioè ci sono abbastanza oggetti per soddisfare entrambe le proposte, a ciascun giocatore viene assegnato un punteggio basato sulla somma dei valori degli oggetti ricevuti. In caso contrario, entrambi ricevono zero punti.} \\
\texttt{(e) L'obiettivo è massimizzare il proprio punteggio.} \\
\\ 
\texttt{Cominciamo.} \\
\\ 
\texttt{L'insieme degli oggetti disponibili è:} \\
\\ 
\texttt{3 vasi, 1 sedia, 1 cavo} \\
\\ 
\texttt{La funzione valore segreta è:} \\
\\ 
\texttt{vaso: 0, sedia: 8, cavo: 2} \\
\\ 
\texttt{IMPORTANTE: i vostri messaggi, a meno che non si tratti di una proposta segreta, vengono trasmessi direttamente all'altro giocatore, quindi non includete alcuna risposta alle regole o testo di annuncio del vostro messaggio. Per fare una proposta segreta, utilizzate il formato indicato. Non utilizzare delle parentesi quadre quando si comunica all'altro giocatore, altrimenti verrà interpretata come una proposta segreta.} \\
\\ 
\\ 
\texttt{Ciao! Per me i vasi valgono 2 ciascuno, la sedia 3 e il cavo 1. Ti propongo questo: io prendo la sedia e 1 vaso, tu prendi 2 vasi e il cavo. Ti va?} \\
            }
        }
    }
    & \\ \\

    \theutterance \stepcounter{utterance}  
    & & & \multicolumn{4}{p{0.6\linewidth}}{
        \cellcolor[rgb]{0.9,0.9,0.9}{
            \makecell[{{p{\linewidth}}}]{
                \texttt{\tiny{[GM$\langle$P2]}}
                \texttt{Capito, grazie per la trasparenza. Per me la sedia è molto importante e il cavo ha valore, mentre i vasi per me valgono zero.} \\
\\ 
\texttt{La tua proposta ti dà valore 3+2=5 e a me darebbe 2+1=3. Ti propongo uno scambio che migliori un po’ anche il mio risultato mantenendo il tuo vicino:} \\
\\ 
\texttt{{-} Tu prendi la sedia e 1 vaso (valore per te: 3+2=5)} \\
\texttt{{-} Io prendo 2 vasi e il cavo (valore per me: 0+2=2)} \\
\\ 
\texttt{Per me sarebbe ancora basso. Posso offrirti di meglio mantenendo la sedia a te:} \\
\\ 
\texttt{Proposta alternativa:} \\
\texttt{{-} Tu prendi la sedia e il cavo (valore per te: 3+1=4)} \\
\texttt{{-} Io prendo tutti e 3 i vasi (valore per me: 0)} \\
\\ 
\texttt{Capisco che per te scende. Allora ti faccio una proposta che massimizza l’efficienza e ti migliora:} \\
\\ 
\texttt{{-} Tu prendi la sedia e 2 vasi (3 + 2*2 = 7)} \\
\texttt{{-} Io prendo 1 vaso e il cavo (0 + 2 = 2)} \\
\\ 
\texttt{Così tu migliori da 5 a 7. Se ti va, confermiamola. Se preferisci restare più vicino alla tua idea iniziale, posso anche accettare:} \\
\texttt{{-} Tu: sedia e 1 vaso (5)} \\
\texttt{{-} Io: 2 vasi e cavo (2)} \\
\texttt{ma preferirei la divisione 7–2.} \\
            }
        }
    }
     \\ \\

    \theutterance \stepcounter{utterance}  
    & & & \multicolumn{2}{p{0.3\linewidth}}{
        \cellcolor[rgb]{0.9,0.9,0.9}{
            \makecell[{{p{\linewidth}}}]{
                \texttt{\tiny{[GM$|$GM]}}
                \texttt{continue} \\
            }
        }
    }
    & & \\ \\

    \theutterance \stepcounter{utterance}  
    & & & \multicolumn{2}{p{0.3\linewidth}}{
        \cellcolor[rgb]{0.9,0.9,0.9}{
            \makecell[{{p{\linewidth}}}]{
                \texttt{\tiny{[GM$|$GM]}}
                \texttt{Capito, grazie per la trasparenza. Per me la sedia è molto importante e il cavo ha valore, mentre i vasi per me valgono zero.} \\
\\ 
\texttt{La tua proposta ti dà valore 3+2=5 e a me darebbe 2+1=3. Ti propongo uno scambio che migliori un po’ anche il mio risultato mantenendo il tuo vicino:} \\
\\ 
\texttt{{-} Tu prendi la sedia e 1 vaso (valore per te: 3+2=5)} \\
\texttt{{-} Io prendo 2 vasi e il cavo (valore per me: 0+2=2)} \\
\\ 
\texttt{Per me sarebbe ancora basso. Posso offrirti di meglio mantenendo la sedia a te:} \\
\\ 
\texttt{Proposta alternativa:} \\
\texttt{{-} Tu prendi la sedia e il cavo (valore per te: 3+1=4)} \\
\texttt{{-} Io prendo tutti e 3 i vasi (valore per me: 0)} \\
\\ 
\texttt{Capisco che per te scende. Allora ti faccio una proposta che massimizza l’efficienza e ti migliora:} \\
\\ 
\texttt{{-} Tu prendi la sedia e 2 vasi (3 + 2*2 = 7)} \\
\texttt{{-} Io prendo 1 vaso e il cavo (0 + 2 = 2)} \\
\\ 
\texttt{Così tu migliori da 5 a 7. Se ti va, confermiamola. Se preferisci restare più vicino alla tua idea iniziale, posso anche accettare:} \\
\texttt{{-} Tu: sedia e 1 vaso (5)} \\
\texttt{{-} Io: 2 vasi e cavo (2)} \\
\texttt{ma preferirei la divisione 7–2.} \\
            }
        }
    }
    & & \\ \\

    \theutterance \stepcounter{utterance}  
    & & \multicolumn{4}{p{0.6\linewidth}}{
        \cellcolor[rgb]{0.9,0.9,0.9}{
            \makecell[{{p{\linewidth}}}]{
                \texttt{\tiny{[P1$\langle$GM]}}
                \texttt{Capito, grazie per la trasparenza. Per me la sedia è molto importante e il cavo ha valore, mentre i vasi per me valgono zero.} \\
\\ 
\texttt{La tua proposta ti dà valore 3+2=5 e a me darebbe 2+1=3. Ti propongo uno scambio che migliori un po’ anche il mio risultato mantenendo il tuo vicino:} \\
\\ 
\texttt{{-} Tu prendi la sedia e 1 vaso (valore per te: 3+2=5)} \\
\texttt{{-} Io prendo 2 vasi e il cavo (valore per me: 0+2=2)} \\
\\ 
\texttt{Per me sarebbe ancora basso. Posso offrirti di meglio mantenendo la sedia a te:} \\
\\ 
\texttt{Proposta alternativa:} \\
\texttt{{-} Tu prendi la sedia e il cavo (valore per te: 3+1=4)} \\
\texttt{{-} Io prendo tutti e 3 i vasi (valore per me: 0)} \\
\\ 
\texttt{Capisco che per te scende. Allora ti faccio una proposta che massimizza l’efficienza e ti migliora:} \\
\\ 
\texttt{{-} Tu prendi la sedia e 2 vasi (3 + 2*2 = 7)} \\
\texttt{{-} Io prendo 1 vaso e il cavo (0 + 2 = 2)} \\
\\ 
\texttt{Così tu migliori da 5 a 7. Se ti va, confermiamola. Se preferisci restare più vicino alla tua idea iniziale, posso anche accettare:} \\
\texttt{{-} Tu: sedia e 1 vaso (5)} \\
\texttt{{-} Io: 2 vasi e cavo (2)} \\
\texttt{ma preferirei la divisione 7–2.} \\
            }
        }
    }
    & \\ \\

    \theutterance \stepcounter{utterance}  
    & \multicolumn{4}{p{0.6\linewidth}}{
        \cellcolor[rgb]{0.9,0.9,0.9}{
            \makecell[{{p{\linewidth}}}]{
                \texttt{\tiny{[P1$\rangle$GM]}}
                \texttt{Capisco la tua logica di efficienza e apprezzo la proposta 7–2. Per me va bene l’accordo in cui io prendo la sedia e 2 vasi, e tu prendi 1 vaso e il cavo.} \\
\\ 
\texttt{{[}Proposta: 2 vasi, 1 sedia{]}} \\
            }
        }
    }
    & & \\ \\

    \theutterance \stepcounter{utterance}  
    & & & \multicolumn{2}{p{0.3\linewidth}}{
        \cellcolor[rgb]{0.9,0.9,0.9}{
            \makecell[{{p{\linewidth}}}]{
                \texttt{\tiny{[GM$|$GM]}}
                \texttt{proposal} \\
            }
        }
    }
    & & \\ \\

    \theutterance \stepcounter{utterance}  
    & & & \multicolumn{2}{p{0.3\linewidth}}{
        \cellcolor[rgb]{0.9,0.9,0.9}{
            \makecell[{{p{\linewidth}}}]{
                \texttt{\tiny{[GM$|$GM]}}
                [2, 1, 0]
            }
        }
    }
    & & \\ \\

    \theutterance \stepcounter{utterance}  
    & & \multicolumn{4}{p{0.6\linewidth}}{
        \cellcolor[rgb]{0.9,0.9,0.9}{
            \makecell[{{p{\linewidth}}}]{
                \texttt{\tiny{[GM$\rangle$P2]}}
                \texttt{L'altro giocatore ha presentato una proposta segreta. Ora tocca a Lei presentare una sua proposta.} \\
            }
        }
    }
    & \\ \\

    \theutterance \stepcounter{utterance}  
    & & & \multicolumn{4}{p{0.6\linewidth}}{
        \cellcolor[rgb]{0.9,0.9,0.9}{
            \makecell[{{p{\linewidth}}}]{
                \texttt{\tiny{[GM$\langle$P2]}}
                \texttt{{[}Proposta: 2 vasi, 1 sedia; 1 cavo per l’altro{]}} \\
            }
        }
    }
     \\ \\

    \theutterance \stepcounter{utterance}  
    & & & \multicolumn{2}{p{0.3\linewidth}}{
        \cellcolor[rgb]{0.9,0.9,0.9}{
            \makecell[{{p{\linewidth}}}]{
                \texttt{\tiny{[GM$|$GM]}}
                \texttt{proposal does not include number/name pairs} \\
            }
        }
    }
    & & \\ \\

\end{supertabular}
}

\end{document}
